% Short research report (Introduction, Methodology, Results, Discussion): what the project found on variance reduction
% in the classical (Euclidean) setting (SAGA), how its estimators and tools were carried over to Muon-type
% (polar-factor) steps, and the two theorems this led to; with an explicit account of what was published before and
% what the project adds (Table 1). Published results are reduced to citations.
% Sources: report/main.pdf (main report), muon/findings.md + muon/results/digest.md (experiments on Muon),
% muon/notes_spectral_theory.pdf (theory, proofs).
% Build (from muon/):  pdflatex report_spectral_muon_new && bibtex report_spectral_muon_new && pdflatex report_spectral_muon_new && pdflatex report_spectral_muon_new
% Figures: muon/report_figures/make_figR1_classical.py, muon/report_figures/make_figR2_muon_vr.py,
%          muon/theory_checks/make_paper_figures.py T2s.
% Bibliography: report_spectral_muon.bib, generated by _scratch/muon_checks/make_report_bib.py muon/report_spectral_muon_new.tex
\documentclass[10pt]{article}
\usepackage[utf8]{inputenc}
\usepackage[T1]{fontenc}
\usepackage{amsmath,amssymb,amsthm,mathtools}
\usepackage{graphicx}
\usepackage{booktabs}
\usepackage{tabularx}
\usepackage{array}
\usepackage{enumitem}
\usepackage{multicol}
\usepackage{etoolbox}
\usepackage[margin=1in]{geometry}
\usepackage[colorlinks=true,linkcolor=blue!55!black,citecolor=blue!55!black,urlcolor=blue!55!black]{hyperref}
\usepackage{xcolor}
\graphicspath{{figures/}}
\setlength{\emergencystretch}{2em}
\clubpenalty=10000 \widowpenalty=10000
\setlength{\textfloatsep}{10pt plus 2pt minus 2pt}
\setlength{\abovecaptionskip}{4pt}
\renewcommand{\topfraction}{0.9}
\renewcommand{\bottomfraction}{0.6}
\renewcommand{\textfraction}{0.08}
\renewcommand{\floatpagefraction}{0.8}
\setcounter{topnumber}{2}
\setcounter{totalnumber}{3}
% compact headings
\makeatletter
\renewcommand\section{\@startsection{section}{1}{\z@}{-2.4ex \@plus -0.6ex \@minus -.2ex}{1.2ex \@plus .2ex}{\normalfont\large\bfseries}}
\renewcommand\subsection{\@startsection{subsection}{2}{\z@}{-1.8ex\@plus -0.5ex \@minus -.2ex}{0.8ex \@plus .2ex}{\normalfont\normalsize\bfseries}}
\makeatother

\newtheorem{theorem}{Theorem}
\newtheorem{proposition}[theorem]{Proposition}

\newcommand{\R}{\mathbb{R}}
\newcommand{\E}{\mathbb{E}}
\newcommand{\Prob}{\mathbb{P}}
\newcommand{\norm}[1]{\left\|#1\right\|}
\newcommand{\nop}[1]{\left\|#1\right\|_{\mathrm{op}}}
\newcommand{\nnuc}[1]{\left\|#1\right\|_{*}}
\newcommand{\nF}[1]{\left\|#1\right\|_{F}}
\newcommand{\inner}[2]{\left\langle #1, #2 \right\rangle}
\newcommand{\abs}[1]{\left|#1\right|}
\newcommand{\tr}{\operatorname{tr}}
\newcommand{\diag}{\operatorname{diag}}
\newcommand{\polar}{\operatorname{polar}}
\newcommand{\Var}{\operatorname{Var}}
\newcommand{\Med}{\operatorname{Med}_\ast}
\newcommand{\Lop}{L_{\mathrm{op}}}
\newcommand{\muop}{\mu_{\mathrm{op}}}
\newcommand{\kop}{\kappa_{\mathrm{op}}}
\newcommand{\kF}{\kappa_{F}}
\newcommand{\lmin}{\lambda_{\min}}
\newcommand{\lmax}{\lambda_{\max}}
\newcommand{\ones}{\mathbf{1}}
\newcommand{\etas}{\eta_\ast}
\newcommand{\muEff}{\mu_{\mathrm{eff}}}
\newcommand{\LEff}{L_{\mathrm{eff}}}
\newcommand{\finding}[1]{\textbf{(#1)}}
\newcommand{\nat}[1]{\textsc{[#1]}}
\newcolumntype{Y}{>{\raggedright\arraybackslash}X}

% compact two-column bibliography (the heading is set by the multicols environment)
\patchcmd{\thebibliography}{\section*{\refname}}{}{}{\typeout{PATCH OF THEBIBLIOGRAPHY FAILED}}
\let\oldthebibliography\thebibliography
\renewcommand{\thebibliography}[1]{\oldthebibliography{#1}\setlength{\itemsep}{0.5pt}\setlength{\parskip}{0pt}}

\title{\vspace{-2.2em}Variance reduction, from SAGA to Muon:\\ where proofs lose, what transfers to polar steps, what limits them}
\author{Project ``Accelerating Variance Reduction'' and its side study on Muon\\ Statistical Learning Theory, UC Berkeley}
\date{\vspace{-0.4em}17 November 2026}

\begin{document}
\maketitle

\begin{abstract}
SAGA minimizes an average of $n$ smooth functions with condition number $\kappa$ in
$O((n+\kappa)\log(1/\varepsilon))$ component gradients; Muon steps along the polar factor of a matrix gradient (or of
its momentum). We ask where the textbook proof of SAGA loses, how far its finite-sum variance reduction carries over to
Muon-type steps, and what limits such steps, and we separate what was published from what we add.
\emph{Methods.} Each inequality of the proof is evaluated exactly along runs by enumerating the
sampled indices; the same enumeration gives the exact mean of polar steps, which are compared on ridge softmax and
least-squares problems at tuned steps on fresh seeds; two theorems are proved and checked numerically. \emph{Results.} On
logistic regression the exact contraction exponent of SAGA is $3.2$ to $220$ times the certified one, a gap that grows
with $\kappa$ through the ratio, up to $81$, between the curvature met along the run and the global constant; yet the
worst-case order cannot be improved by tuning the proof, by a better analysis of constant-step SAGA, or by a fixed
preconditioner. Inside Muon-type steps the same estimators act on the fixed point and not only on the variance; at
batch size one the nuclear-scaled polar step of SAGA behaves as the MISO algorithm; and known Euclidean devices match
the gain of orthogonalization over SAGA. For ridge-regularized generalized linear models whose loss curvature vanishes
along a sequence of parameters, no norm has a better worst-case condition number than the Euclidean one; with
additive gradient noise and suitably decreasing step lengths, stochastic polar steps on smooth strongly convex objectives converge almost surely to a point
where the gradient is the negative of a \emph{nuclear median} of the noise, whereas a finite-sum estimator has the
minimizer as its only fixed point. We propose no method that beats the state of the art.
\end{abstract}

\section{Introduction}\label{sec:intro}

\paragraph{Context.} Variance-reduced methods such as SVRG \cite{johnson2013accelerating} and SAGA
\cite{defazio2014saga} build, for a finite sum $F=\frac1n\sum_if_i$, a gradient estimator whose error vanishes at the
minimizer, and converge linearly with a constant step; Proposition~5.9 of \cite{bach2024learning} proves it for SAGA,
and Katyusha \cite{allenzhu2018katyusha,kovalev2020dont} attains the optimal order
$(n+\sqrt{n\kappa})\log(1/\varepsilon)$ \cite{lan2018optimal}. Muon \cite{jordan2024muon} steps along an
approximation (by Newton--Schulz iterations) of the polar factor of the momentum of the gradient of each weight matrix; without momentum this is
steepest descent for the spectral norm \cite{carlson2015stochastic,bernstein2024old}. Its theory is recent: the
dynamics on least squares \cite{shulgin2026quadratic,vasudeva2025muon} and the effect of gradient noise
\cite{xie2026muon,chen2026when,paquette2026phases} are being worked out, and its published variance-reduced versions
are mostly momentum-based \cite{liu2025marsm}.

\paragraph{Questions and approach.} The main project asked where the proof of Proposition~5.9 is loose and what
acceleration needs. We treat the proof as a measurable object: each of its inequalities is evaluated exactly along
trajectories, so that the gap between the certified and the actual contraction can be attributed to them. The answer
has two halves (Section~\ref{sec:classical}): in the worst case the theorem has the right order, and neither a better
analysis nor, on logistic regression, a fixed metric improves it; but along the trajectories followed the curvature is far
above the global constant, and this is where the proof loses. A side study then asked to what extent the finite-sum
estimators of the main project apply to Muon-type steps. Muon is a change of geometry used with stochastic gradients, so two questions arise:
can the spectral geometry improve the worst case, which no fixed metric could, or only the trajectories, and what
does an estimator designed for steps that are linear in it do inside a step that is not?

\paragraph{Contributions.} Table~\ref{tab:sota} sets each result against the published state of the art: a way of
measuring a proof along trajectories, what it shows for SAGA and loopless Katyusha, and what cannot be gained
\finding{C1}--\finding{C6}; an impossibility result for every norm on ridge logistic and softmax models
(Theorem~\ref{thm:nfl}); the limit of stochastic polar steps with additive noise and decreasing lengths, a point where
the gradient is minus a nuclear median of the noise (Theorem~\ref{thm:smd}), and the fixed point without offset of
finite-sum estimators (Proposition~\ref{prop:comp}(c)); and, for variance reduction inside Muon-type steps, the
reading of the nuclear-scaled polar step of SAGA as MISO, a variant with a certified rate and a benchmark whose
outcome is negative \finding{M1}--\finding{M4}. Section~\ref{sec:method} describes the methods,
Section~\ref{sec:results} the results and Section~\ref{sec:discussion} what transferred from SAGA to Muon and why.
Proofs and all numbers are in the main report (\texttt{report/main.pdf}), the account of the experiments on Muon
(\texttt{muon/findings.md}, tables in \texttt{muon/results/digest.md}) and the theory note
(\texttt{muon/notes\_spectral\_theory.pdf}).

\begin{table}[t]
\caption{What was published and what this study adds. Tags (Section~\ref{sec:status}): \nat{p} proved, \nat{m}
measured exactly along runs, \nat{e} experiment, \nat{f} formal expansion checked by simulation. The additions were not
found by the literature checks of Section~\ref{sec:status}, with the exceptions listed there.}\label{tab:sota}
\centering\footnotesize
\renewcommand{\arraystretch}{1.1}
\begin{tabularx}{\textwidth}{@{}>{\raggedright\arraybackslash}p{0.15\textwidth}YY@{}}
\toprule
Question & Published state of the art & Added by this study \\
\midrule
How loose is the proof of SAGA, and where? \finding{C1}--\finding{C3} &
Rate and Lyapunov proof \cite{defazio2014saga,bach2024learning}; worst case over a class of functions by performance
estimation and matrix inequalities \cite{taylor2017smooth,hu2017unified}; one-point curvature conditions
\cite{zhang2013gradient}. &
Each inequality of the proof evaluated exactly along runs \nat{m}: exact exponent $3.2$--$220$ times the certified
one; $77$--$96\%$ of the gap due to the two global constants ($\kappa\ge10^4$); curvature along the run up to
$81\mu$. \\
What can tuning or a better analysis gain? \finding{C4} &
A sharper step and rate \cite{hofmann2015variance}; order $\kappa\log\frac1\varepsilon$ for constant steps of one-step
methods \cite{arjevani2016lower} and for steps that depend only on the smoothness constant \cite{arjevani2016iteration}; lower bounds for classes of methods
\cite{lan2018optimal,woodworth2016tight}. &
These are the best constants of the textbook proof, within a factor $3$ of its own \nat{p}; an elementary instance
where every constant step needs $\frac{\kappa-1}4\log\frac1{2\varepsilon}$ iterations \nat{p}; a memory ceiling of
$1/n$ per iteration for this Lyapunov function \nat{p}. \\
Can a change of geometry improve the worst case? \finding{C5}, Thm.~\ref{thm:nfl} &
SAGA in a fixed metric \cite{dragomir2021fast,kolte2015accelerating}; preconditioning helps where the loss curvature
is bounded below (square loss; logistic loss on bounded sets) \cite{yang2014data,gonen2016solving}, and no linear
preconditioner helps some functions \cite{hird2025quantifying}; Muon as spectral steepest descent
\cite{bernstein2024old}. &
When the loss curvature vanishes along a sequence (logistic, softmax), no fixed preconditioner and no norm improves
the global condition number or the per-sample ones \nat{p}. \\
What does acceleration need? \finding{C6} &
Katyusha, whose negative momentum cancels an inner-product term of the variance bound \cite{allenzhu2018katyusha},
and its loopless version \cite{kovalev2020dont}; without it, momentum accelerates with minibatches of order
$n^{2/3}$ and a worse dependence on $n$ \cite{driggs2022accelerating}. &
In the Lyapunov proof of \cite{kovalev2020dont}, pessimistic by $2.8$--$5.9$ \nat{m}, this cancellation is isolated
at a single step \nat{p}; without it, divergence on hard quadratics but half the gradients on a logistic problem,
consistent with \cite{allenzhu2018katyusha} \nat{e}. \\
What does gradient noise do to a polar step? Thm.~\ref{thm:smd}, Props.~\ref{prop:comp}--\ref{prop:offset} &
Noise smooths the polar map \cite{xie2026muon,chen2026when}; sign methods follow a median and converge under
unimodal symmetric noise \cite{chen2020distributed,bernstein2019majority}; in dimension one Muon need not converge
under skewed noise \cite{do2026muon}. &
With additive noise, polar steps are stochastic mirror descent on a smoothed nuclear norm and, for decreasing
lengths, converge to a point where the gradient is minus a nuclear median of the noise \nat{p}; at batch one on least
squares ($K,d\ge2$), to a least-distances estimator \nat{p}; for i.i.d.\ skewed entries, a limit gradient of order $1/b$,
per entry a factor of order $\max(K,d)$ below sign descent \nat{f}. \\
What does a finite-sum estimator change? Prop.~\ref{prop:comp}(c), \finding{M1}, \finding{M2}, \finding{M4} &
Momentum-based variance reduction for Muon \cite{liu2025marsm}; SVRG in a fixed-length polar step
\cite{li2026implicit} and in sign descent \cite{chzhen2023signsvrg}; floors of fixed lengths \cite{lau2025polargrad}; MISO \cite{mairal2015incremental}. &
The minimizer becomes the only fixed point \nat{p}; rounding precision needs a step that vanishes with the estimator
\nat{e}; at batch one the nuclear-scaled polar step of SAGA behaves as MISO \nat{e}; a frozen polar metric inherits
the rate of Proposition~5.9 \nat{p}. \\
Is orthogonalization worth it for finite sums? \finding{M3} &
Damped SAGA \cite{driggs2019biased}, importance sampling \cite{xiao2014proximal}, covariance preconditioning
\cite{yang2014data}, for Euclidean steps. &
A tuned comparison on six instances \nat{e}: where conditioning limits, orthogonalization at every step divides the
passes of SAGA by $1.3$--$5$ at batch one, as these Euclidean devices do (ratio $0.78$--$1.11$); where memory limits,
nothing helps. \\
\bottomrule
\end{tabularx}
\end{table}

\section{Methodology}\label{sec:method}

\subsection{Problems and algorithms}\label{sec:setting}

\paragraph{Finite sums and SAGA.} $F=\frac1n\sum_{i=1}^nf_i$ on $\R^d$, each $f_i$ convex and $R^2$-smooth, $F$
$\mu$-strongly convex with minimizer $\etas$, $\kappa=R^2/\mu$; costs are counted in component gradients. SAGA keeps a
table $z_i$ of past component gradients and steps along $f_i'(\theta)-z_i+\bar z$, an unbiased estimator whose error
vanishes at $\etas$. Proposition~5.9 of \cite{bach2024learning}: with the step $1/(4R^2)$ a Lyapunov function (squared
distance plus a multiple of the memory error) contracts in expectation by $1-\min\{1/(3n),3\mu/(16R^2)\}$ per
iteration.

\paragraph{Matrix problems and polar steps.} For $W\in\R^{K\times d}$, $\nF\cdot$, $\nop\cdot$, $\nnuc\cdot$ are the
Frobenius, spectral and nuclear norms, and $\polar(\omega)=UV^\top$ for a singular value decomposition
$\omega=U\diag(s)V^\top$ with $s>0$ (so $\polar(0)=0$). A \emph{Muon-type step} is $W^+=W-\eta\,\polar(\omega)$, where $\omega$ estimates
$\nabla F(W)$; its length is fixed, as in Muon, or \emph{nuclear-scaled}, $\eta=\gamma\nnuc\omega$. For $F\in C^2$
with minimizer $W_\ast$ and a norm $N$, $\mu_N,L_N$ are the best constants in
$\mu_NN(\Delta)^2\le D^2F(W)[\Delta,\Delta]\le L_NN(\Delta)^2$ and $\kappa_N=L_N/\mu_N$ (subscripts ``op'' and $F$
for the spectral and Frobenius norms). The models are ridge-regularized generalized linear models, (GLM):
$F(W)=\frac1n\sum_i\ell_i(Wx_i)+\frac\lambda2\nF W^2$ with $x_i\in\R^d$, $\lambda>0$ and convex $C^2$ losses,
$c_i=\sup_z\lmax(\nabla^2\ell_i(z))<\infty$ and $f_i(W)=\ell_i(Wx_i)+\frac\lambda2\nF W^2$; softmax regression and
multi-output least squares are the instances used.

\subsection{Measuring a proof along trajectories}\label{sec:measure}

We rewrite the proof of Proposition~5.9 as an exact one-step identity for the conditional expectation of the Lyapunov
function, followed by six labelled inequalities: square norm (I1), cocoercivity with $R^2$ (I2), variance
(I3), strong convexity with $\mu$ (I4), minimum of two rates (I5) and initialization (I6, which only enters a
constant). The sampled index being uniform over $n$ values, every
conditional expectation is computed exactly along a run by enumerating the indices: this gives the exact one-step
contraction and the slack of each inequality, and, because undoing I1--I5 one at a time gives a monotone chain from
the certified to the exact contraction, an additive attribution of the logarithmic gap between the two (inequalities
undone in reverse proof order). We also record the curvature $\muEff(\theta)=\inner{\theta-\etas}{F'(\theta)}/\norm{\theta-\etas}^2$,
the restricted-secant constant \cite{zhang2013gradient} at the iterate, and its cocoercivity counterpart
$\LEff(\theta)=\frac1n\sum_i\norm{f_i'(\theta)-f_i'(\etas)}^2/\inner{\theta-\etas}{F'(\theta)}$, which
turn I4 and I2 into identities. Performance estimation and matrix inequalities compute the worst case over a class of
functions \cite{taylor2017smooth,hu2017unified}; this measures how far each inequality of a given proof is from
equality along one trajectory. Data: ridge logistic regression with $n=1000$, $d=50$, $\kappa$ from $10$ to $10^5$,
covariance spectra $j^{-\alpha}$ with $\alpha\in\{0,1,2\}$, five sampling seeds, exponents averaged over passes 2--20.
The same instrumentation is applied to a Lyapunov proof of loopless Katyusha written in the same style (logistic regression with
$\kappa\in\{10^3,10^4\}$, $\alpha=1$, and one hard quadratic), and gradient counts are measured on four quadratic
instances with $n=16$ (``hard quadratics''), one of them the lower-bound instance of \cite{lan2018optimal}.

\subsection{Variance reduction inside Muon-type steps}\label{sec:vrmethod}

The estimators of SAGA, loopless SVRG and loopless Katyusha are written $\omega=m+\vartheta c$: a predictable part
$m$ (memory mean plus the gradient of the regularizer) and an innovation $c$, of rank one at batch size one;
$\vartheta=1$ for SAGA, $\vartheta<1$ damps the innovation, and the Euclidean step $\gamma(m+\vartheta c)$ with
$\vartheta=1/n$, $\gamma=1/\lambda$ is MISO (minimization by incremental surrogate optimization) with lower-bound
surrogates \cite{mairal2015incremental}. A method chooses the estimator, where the orthogonalization acts (on
$\omega$: the ``joint'' SAGA-Muon; on $m$ only; or through a \emph{polar metric} $\omega\mapsto\omega C$ read off the
polar decomposition of a gradient, frozen at $W_0$ (``one-shot'') or refreshed along the run), the length (fixed,
nuclear-scaled, or clipped, $\min\{\rho,\gamma\nnuc\omega\}$) and the orthogonalization (exact, Newton--Schulz, or the
bounded-gain ``soft'' map $U\diag(s/\sqrt{s^2+\epsilon^2})V^\top$, $\epsilon$ proportional to $\nF{\vartheta c}$). The
Euclidean baselines are SAGA, SAGA with a damped innovation, with importance sampling, and in the feature-covariance
metric. Floors and passes are measured on runs; the one-step expectations behind \finding{M2} are computed exactly by
enumerating the $n$ indices, as in Section~\ref{sec:measure}. Of six instances ($n\in\{200,500,1000\}$,
$K\times d\in\{8\times40,10\times50,20\times20\}$), two are memory-limited ($R^2/\mu$ of the order of $n$) and four
conditioning-limited ($R^2/\mu$ from $9n$ to $50n$; two of them square). Costs are in passes over the data
(factorizations not counted); the target is a relative squared distance $10^{-6}$ to the exact minimizer, $10^{-9}$ in
the tuned runs on the memory-limited instances. Each family (a method with its grid of step constants) is tuned on
five seeds, then run on ten fresh ones; ten families are re-tuned and rerun on two further draws of the data for the
two instances of Figure~\ref{fig:muon}(c). Differences below about $10\%$ are not resolved.

\subsection{Analysis of stochastic polar steps}\label{sec:noisemethod}

Let $\xi$ be a random $K\times d$ matrix (the gradient noise), $\psi_\xi(G)=\E[\nnuc{G+\xi}-\nnuc\xi]$ and
$\Phi_\xi(G)=\E\,\polar(G+\xi)$, the \emph{mean field} of a stochastic polar step at the gradient $G$. No moment is
required; $\psi_\xi$ is convex and $\Phi_\xi(G)\in\partial\psi_\xi(G)$, the randomized-smoothing identity
\cite{duchi2012randomized,jiang2026adaptive}. We call $\Med(\xi)=\arg\min_C\E[\nnuc{\xi-C}-\nnuc\xi]$ the set of
\emph{nuclear medians} of $\xi$ (the median of a law in a normed space \cite{minsker2015geometric}). The analysis
reads a stochastic polar step as stochastic mirror descent on $\psi_\xi$ in the space of gradients, with mirror map
$F^\ast$, as in \cite{maddison2021dual} for smooth maps of the gradient (here the nonsmooth map $\polar$); it concludes
with the Robbins--Siegmund lemma and the dual-norm Polyak--\L{}ojasiewicz inequality and its converse
\cite{balles2020geometry,leplat2025norm}, and sizes the offset of the limit for minibatches by a one-term Edgeworth
expansion. Predictions are checked numerically: zeros of the mean field from exactly sampled batch means (enumerated at batch one), constants by quadrature or Monte Carlo, and runs on an unregularized finite-sum least-squares problem with an intercept and
skewed residuals.

\subsection{Status of the claims}\label{sec:status}

Each addition is tagged \nat{p} proved (proofs in the main report, the theory note or \texttt{muon/notes\_theory.md}),
\nat{m} measured exactly along runs, \nat{e} experiment, or \nat{f} formal expansion checked by simulation. ``Not
found'' refers, for Section~\ref{sec:classical}, to the literature review of the main project (not a separate search
for each statement; targeted automated searches of 17 November 2026 for the method of \finding{C1}--\finding{C2}, the lower
bound of \finding{C4}, \finding{C5} and \finding{C6}) and, for the rest, to automated literature checks of 12 and 16
November 2026 that opened the closest sources (about sixty papers for the theory, several of them a few weeks old).
Written after these checks and searched only in an earlier form: Theorem~\ref{thm:smd} (the mean-field iteration), the
constant of \eqref{eq:offset} (the order $1/b$) and the uniqueness and convergence clauses of
Proposition~\ref{prop:comp} (stationary and fixed points); not searched: Theorem~\ref{thm:nfl} for an arbitrary norm
(the checks covered right-preconditioned spectral norms).

\section{Results}\label{sec:results}

\subsection{The classical case: where the proof of SAGA loses}\label{sec:classical}

\begin{figure}[t]
\centering
\includegraphics[width=0.97\textwidth]{figR1_classical}
\caption{SAGA with the textbook step on ridge logistic regression ($n=1000$, $d=50$). (a) $\kappa=10^4$, spectrum
$j^{-1}$: certified bound of Proposition~5.9, a-posteriori bound with the effective constants of \finding{C3} along
one run, and squared distance to the minimizer (mean of five runs). (b) Exact over certified contraction exponent
(passes 2--20, five seeds) against $\kappa$, for spectra $j^{-\alpha}$. (c) Share of the logarithmic gap carried by
each inequality ($\alpha=1$; undone in reverse proof order; mean and range over five seeds).}\label{fig:classical}
\end{figure}

\begin{description}[leftmargin=2.3em,labelwidth=2.0em,labelsep=0.3em,itemsep=1pt,topsep=3pt,parsep=0pt,font=\normalfont\bfseries]
\item[(C1)] \emph{The gap.} The exact contraction exponent is $3.2$--$3.7$ times the certified one for
$\kappa\le10^2$ and $19$--$220$ times at $\kappa=10^5$ (Figure~\ref{fig:classical}(b)).

\item[(C2)] \emph{Its anatomy.} For $\kappa\ge10^4$ I2 and I4, which replace a quantity depending on the iterate by a
global constant, carry $77$--$96\%$ of
the logarithmic gap; for $\kappa\le10^2$, where the memory term is active, I1 ($39$--$44\%$) and
I2 ($32$--$37\%$) carry it (Figure~\ref{fig:classical}(c)). The variance step I3 is nearly tight (at most $0.6\%$ of the
gap). The part that grows with $\kappa$ is the ratio $\muEff/\mu$, up to $81$: I1 and I2 together contribute at most a
factor $8/3$ ($2.45$--$2.62$ measured), and the product of the two factors reproduces $91$--$96\%$ of the speed-ups
for $\kappa\ge10^4$.

\item[(C3)] \emph{Effective constants explain, they do not certify.} The global strong-convexity constant of ridge
logistic regression is exactly $\lambda$, whatever the data, whereas near the optimum $\muEff$ is at least about the
smallest eigenvalue of the Hessian at $\etas$, a value that SAGA reaches when its error aligns with the flattest
direction ($\muEff/\mu=9.94$ against $9.86$, $\kappa=10^4$, $\alpha=0$, one run). Rewritten with $\muEff$ and
$\LEff$, the bound of Proposition~5.9 recovers $59$--$96\%$ of the logarithmic gap for $\kappa\ge10^4$
(Figure~\ref{fig:classical}(a)): an a-posteriori reading, not a guarantee.

\item[(C4)] \emph{What cannot be gained.} The best step and Lyapunov weight within the proof gain at most a factor
$3$, and at most $4/3$ when $\kappa\ge9n/16$ (the optimum is the step and rate of \cite{hofmann2015variance}),
although SAGA still converges with a step eight times the supremum of the steps the proof admits ($\kappa=10^3$, $\alpha=1$,
three seeds). No analysis removes the
exponent $1$ of $\kappa$: on a two-dimensional quadratic with $n$ identical components, SAGA with any constant step,
even one that depends on $\mu$, needs at least $\frac{\kappa-1}4\log\frac1{2\varepsilon}$ iterations (the mean iterate is
that of gradient descent, so the order $\kappa$ is that of one-step stationary methods \cite[Theorem~8]{arjevani2016lower};
for steps that depend only on $R^2$ it also follows from \cite{arjevani2016iteration}). And no bound on the expectation of
this Lyapunov function certifies more than $1/n$ per iteration, since table entries not yet refreshed keep their
error: a \emph{memory ceiling} of about $e^{-1}$ per pass.

\item[(C5)] \emph{No free lunch for preconditioning.} Proposition~5.9 holds in any fixed quadratic metric $P$ with the
relative constants $R_P^2=\max_iL_P(f_i)$ and $\mu_P=\mu_P(F)$ of the norm $\norm{P^{1/2}\cdot}$ (as for Bregman-SAGA
\cite{dragomir2021fast}), and a metric built from the data covariance improves them when the loss curvature is bounded
below where the predictions lie, as for ridge regression \cite{yang2014data,gonen2016solving} and, on bounded sets,
logistic regression \cite{yang2014data}. For ridge logistic regression, whose loss curvature vanishes at infinity,
every $P$ has $\kappa_P=R_P^2/\mu_P\ge\kappa$ for these global constants. Yet at $\kappa=10^3$ the statistical metric $P=\hat\Sigma+\lambda I$ ($\hat\Sigma$ the
feature covariance) raises the observed exponent per pass from $1.00$ to $1.11$ ($\alpha=1$) and from $0.61$ to
$1.26$ ($\alpha=2$), while the certified bound ranks it below the identity: the gain is a property of the trajectory,
accounted for, heuristically, only by local constants at the optimum.

\item[(C6)] \emph{Acceleration.} In the Lyapunov proof of loopless Katyusha \cite{kovalev2020dont}, the negative
momentum, whose role is to cancel an inner product of the variance bound \cite{allenzhu2018katyusha}, acts at a single
step, an exact cancellation; without it the same argument certifies at best $1-\kappa^{-2}$ per iteration (a limit of
the argument, not a lower bound). The certificate is pessimistic by a factor $2.8$--$5.9$ over passes 2--20 ($5.2$ and $15.7$ for
Proposition~5.9 on the same logistic problems). Without negative momentum the method diverges on hard quadratics for
every $\kappa\ge10^3$, but on a logistic problem ($\kappa=10^4$, one seed) needs about half the gradients of loopless
Katyusha with its worst-case parameters: a worst-case safeguard can cost on benign instances. The exponents $1$ and $\frac12$ of $\kappa$ predicted for SAGA and loopless
Katyusha are observed on hard quadratics (log--log slopes $0.95$--$1.02$ and $0.49$--$0.52$).
\end{description}

\subsection{Variance reduction inside Muon-type steps}\label{sec:transfer}

\begin{figure}[t]
\centering
\includegraphics[width=0.97\textwidth]{figR2_muon_vr}
\caption{Finite-sum variance reduction inside Muon-type steps. (a) Softmax, memory-limited instance, one seed,
constant parameters (SAGA estimator with the fixed length $10^{-2}$ or the nuclear-scaled one; Muon with
$\mathrm{lr}=10^{-2}$ and the minibatch gradient with the nuclear-scaled length, both with $b=32$): relative squared
distance to the minimizer. (b) Same instance ($\kappa=n$), untuned: joint SAGA-Muon and MISO with the same sampled
indices, and SAGA at its textbook step ($c$ in the legend: step constant). (c) Passes to $10^{-6}$ at tuned steps on the two non-square conditioning-limited instances
(median and range over ten seeds; arrow: a seed misses the target). B-SAGA: damped innovation; IS: importance
sampling; cov-metric: feature-covariance metric.}\label{fig:muon}
\end{figure}

\begin{description}[leftmargin=2.3em,labelwidth=2.0em,labelsep=0.3em,itemsep=1pt,topsep=3pt,parsep=0pt,font=\normalfont\bfseries]
\item[(M1)] \emph{Two floors, two ingredients} (Figure~\ref{fig:muon}(a)). Polar steps of fixed length stall at a
floor \cite{lau2025polargrad}. With the SAGA estimator and a fixed length $\eta$, the iterates settle at an
absolute squared distance $3.8\,\eta^2$ (softmax) or $6.0\,\eta^2$ (least squares): variance reduction makes the floor
quadratic in the length (it is linear for minibatch Muon) but does not remove it, and a momentum-based buffer (as in
MARS-M \cite{liu2025marsm}) keeps a noise floor at constant parameters whatever the length. Only a finite-sum
estimator with a step that vanishes with it (nuclear-scaled or clipped length; on least squares also the soft map)
reaches rounding precision.

\item[(M2)] \emph{At batch size one, the polar factor of the SAGA estimator is MISO in disguise}
(Figure~\ref{fig:muon}(b)). The predictable part $m$ is the gradient of an isotropic quadratic, the orthogonalized
matrix is dominated by the rank-one innovation, and the nuclear-scaled step drives $m$ to zero; the method that keeps
$m=0$ exactly is MISO (Section~\ref{sec:vrmethod}). With the same sampled indices the
traces follow each other: the joint SAGA-Muon needs the passes of MISO where MISO converges ($11.8$ against $11.9$,
medians of three seeds, softmax with $\kappa=n$; SAGA at its textbook step: $20.2$) and fails where MISO fails
($\kappa\gg n$). It is not a spectral method.

\item[(M3)] \emph{Head-to-head at tuned steps} (Figure~\ref{fig:muon}(c)). Where the memory limits, the best family
of every variance-reduced class (Euclidean, orthogonalized, or in a fixed metric) needs $16$--$18$ passes (softmax) and
$17$--$21$ (least squares) to $10^{-9}$: every geometry runs near the rate of the memory ceiling of \finding{C4},
$e^{-1}$ per pass. Where the conditioning limits (the two non-square instances: softmax / least squares), SAGA
needs $69.5$ / $98.7$ passes, the joint
SAGA-Muon $29.0$ / $43.5$ (nine seeds of ten in the second case), its soft version $23.9$ / $29.8$, and damped SAGA in
a polar metric frozen once \finding{M4} $22.1$ / $24.4$; SAGA with a damped innovation \cite{driggs2019biased} needs
$33.8$ / $31.9$, with importance sampling \cite{xiao2014proximal} in addition $25.7$ / $27.6$, and in the
feature-covariance metric of \finding{C5} \cite{yang2014data} $18.1$ / $16.6$ ($15.1$ / $16.1$ with damping). Over
three draws of the data, the best method recomputing the polar factor needs $0.20$--$0.37$ of the passes of SAGA and
$0.78$--$1.08$ of those of the best Euclidean method without preconditioner; the covariance metric is the best method
on least squares for every draw, and on softmax for one draw of three (undamped; the damped version was not rerun).

\item[(M4)] \emph{A frozen polar metric makes the proof of SAGA apply.} If the map $\omega\mapsto\omega C$, with
$C\succ0$ proportional to $(D^\top D)^{-1/2}$ on the row space of one full gradient $D=\nabla F(W_0)$ (gains floored)
and to the identity on its complement, is computed once,
the method is SAGA on a reparametrized problem and the proof of Proposition~5.9 goes through in the metric: with the
step $1/(4R_P^2)$ a Lyapunov function contracts by $1-\min\{1/(3n),3\mu_P/(16R_P^2)\}$ per iteration. For least
squares these constants are exact and can be far better than the Euclidean ones ($R_P^2/\mu_P=42$ against
$R^2/\mu=4364$ on the square least-squares instance, where the polar factor of the gradient is close to the Newton direction
\cite{shulgin2026quadratic}); for softmax they never improve (Theorem~\ref{thm:nfl}). On that instance this
``one-shot polar-metric SAGA'' needs $13.3$ passes (SAGA $27.1$, covariance metric $13.1$).
\end{description}

\noindent What does not transfer: orthogonalizing a loopless-SVRG estimator (at most one seed of ten reaches the
target), a spectral step inside loopless Katyusha (same passes as the Euclidean method), and a polar metric refreshed
along the run.

\subsection{Conditioning: \finding{C5} holds for every norm}\label{sec:theory}

\begin{theorem}[no free lunch, for every norm]\label{thm:nfl}
Consider (GLM) and assume that there is a sequence $W_t$ with $\max_i\nop{\nabla^2\ell_i(W_tx_i)}\to0$ (logistic and
softmax losses with nonzero $x_i$: $W_t=tW_1$ for almost every $W_1$). Let $N$ be any norm on $\R^{K\times d}$ and
$c_N=\max_{\nF\Delta=1}N(\Delta)$. Then $\mu_N(F)=\lambda/c_N^2$ and $L_N(g)\ge L_F(g)/c_N^2$ for every convex $C^2$
function $g$. Consequently
$\kappa_N(F)\ge\kF(F)$ and $L_N(f_i)/\mu_N(F)\ge L_F(f_i)/\mu_F(F)=(c_i\norm{x_i}^2+\lambda)/\lambda$ for every $i$.
\end{theorem}
The proof is three lines: the ridge term gives the floor $D^2F\ge\lambda\nF\cdot^2\ge(\lambda/c_N^2)N^2$, attained in
the limit along $W_t$ in a direction maximizing $N/\nF\cdot$, and $N\le c_N\nF\cdot$ bounds every smoothness constant
from below. This covers the spectral norm and its preconditioned versions, the $\ell_\infty$ norm of sign descent and
Kronecker-factored metrics, for the objective and for the per-sample constants of variance-reduced methods; for $K=1$
and Euclidean metrics it is \finding{C5} (functions that no linear preconditioner improves exist in every dimension
\cite[Proposition~7]{hird2025quantifying}). The curvature floor of a ridge model is isotropic and a norm can only trade
smoothness against strong convexity: any advantage of a step in another norm on such problems is a property of the
trajectory. The hypothesis on the losses is needed: for ridge least squares the norm $\nF{\Delta\Sigma^{1/2}}$,
$\Sigma$ the Hessian, has condition number one. The spectral geometry has two further costs: $\kop\ge\min(K,d)$ for
every $F$, the matrix form of a folklore fact on $\ell_\infty$ \cite{sidford2018coordinate}; and for least squares with
$K\ge d$, $\Lop(F)=\tr\Sigma$ \cite{davis2025spectral} is the \emph{average} of the $\Lop(f_i)$, so that without ridge
the full-batch nuclear-scaled step has the condition number of importance-sampled variance reduction
\cite{xiao2014proximal} at $n$ gradients per step (this reading was not found).

\subsection{Noise: nuclear medians set the limit}\label{sec:noise}

With the objects of Section~\ref{sec:noisemethod}, $\arg\min\psi_\xi=-\Med(\xi)$ is nonempty, convex and compact, it
contains every zero of $\Phi_\xi$, and it consists of zeros of $\Phi_\xi$ where $G+\xi$ has full rank almost surely.

\begin{theorem}[stochastic polar steps are stochastic mirror descent]\label{thm:smd}
Let $F$ satisfy $0<\muop\le\Lop<\infty$, let $\xi_t$ be i.i.d.\ copies of $\xi$ independent of the past, $\eta_t\ge0$
deterministic lengths, and $W_{t+1}=W_t-\eta_t\,\polar(\nabla F(W_t)+\xi_t)$, $G_t=\nabla F(W_t)$. For every
$G^\dagger\in\arg\min\psi_\xi$, with $W^\dagger=\nabla F^\ast(G^\dagger)$ and $D_F$ the Bregman divergence of $F$:
\begin{enumerate}[label=(\alph*),itemsep=0pt,topsep=2pt]
\item $\E[D_F(W_{t+1},W^\dagger)\mid W_t]\le D_F(W_t,W^\dagger)-\eta_t(\psi_\xi(G_t)-\min\psi_\xi)+\frac{\Lop}2\eta_t^2$,
hence $\E\,\psi_\xi(\bar G_T)-\min\psi_\xi\le\big(D_F(W_0,W^\dagger)+\frac{\Lop}2\sum_{t<T}\eta_t^2\big)/\sum_{t<T}\eta_t$
for the weighted average $\bar G_T=\sum_{t<T}\eta_tG_t/\sum_{t<T}\eta_t$, whatever the lengths;
\item if $\sum_t\eta_t=\infty$ and $\sum_t\eta_t^2<\infty$, $G_t$ converges almost surely to a point of $-\Med(\xi)$;
\item hence, for such lengths, almost surely: $W_t\to W_\ast$ if $\Med(\xi)=\{0\}$, and if $0\notin\Med(\xi)$ then
$\lim_tF(W_t)-F_\ast\ge\min_{C\in\Med(\xi)}\nnuc C^2/(2\Lop)>0$.
\end{enumerate}
\end{theorem}
The theorem replaces the usual condition $\E\xi=0$ by $\Med(\xi)=\{0\}$; otherwise the limit gradient, a point of
$-\Med(\xi)$, is an \emph{offset}. Its ingredients are published (Section~\ref{sec:noisemethod}), and so are its scalar and
coordinatewise analogues:
sign methods drive a median of the stochastic gradient to zero \cite{chen2020distributed}, and normalized steps on a
location problem estimate the geometric median \cite{cardot2013efficient}. The statement for polar steps was not found.

\begin{figure}[t]
\centering
\includegraphics[width=0.97\textwidth]{figT2s_noise}
\caption{Stochastic polar steps on finite-sum least squares with an intercept and skewed residuals ($n=2000$, $K=2$,
$d=3$). (a) Batch one, lengths $\eta_t=0.3\,(1+t/200)^{-3/4}$ from $W_0=0$: with the minibatch gradient the iterate
averaged over $(t/2,t]$ and over $32$ runs converges to the least-distances solution $W_{\mathrm{LAD}}$ of
Proposition~\ref{prop:comp}(d); with the SAGA estimator, after an excursion while its table (initialized at $W_0$) is
stale, to the least-squares solution $W_{\mathrm{LS}}$ (one pass: $2000$ iterations). (b) Batch size times the nuclear
norm of the gradient at the limit point: zeros of the mean field (line), runs with decreasing lengths (dots; same
averaging, $64$ runs), SAGA estimator with the nuclear-scaled length (squares, at rounding level).}\label{fig:theory}
\end{figure}

\begin{proposition}[when is there an offset?]\label{prop:comp}
(a) If $\xi$ and $-\xi$ have the same law, $0\in\Med(\xi)$, and $\Med(\xi)=\{0\}$ if moreover $0$ is in the support
of $\xi$.
(b) If $\xi$ is integrable and $\E\xi=0$, every $C\in\Med(\xi)$ has $\nnuc C\le\E\nnuc\xi$, so the limit of
Theorem~\ref{thm:smd} has $F(W^\dagger)-F_\ast\le(\E\nnuc\xi)^2/(2\muop)$.
(c) For a SAGA or SVRG estimator, the stochastic gradient at $W_\ast$, with its memory (table or reference point) also
at $W_\ast$, is $0$ for every index: a Dirac mass, whose median is its mean. Such estimators have a fixed point
without offset (no index moves it), and it is the only one.
(d) (Batch one on least squares.) Let $g_j(W)=(Wx_j-y_j)x_j^\top$, $x_j\ne0$, $j$ uniform in $\{1,\dots,n\}$. The
batch-one polar step is the stochastic subgradient method on $\Lambda(W)=\frac1n\sum_i\norm{Wx_i-y_i}/\norm{x_i}$;
with deterministic lengths, $\sum_t\eta_t=\infty$, $\sum_t\eta_t^2<\infty$, it converges almost surely to a minimizer
of this weighted
least-distances objective (least absolute deviations for $K=1$), in general not the least-squares solution.
\end{proposition}
Item (a) parallels the unimodal symmetric-noise condition under which signSGD converges
\cite{bernstein2019majority}; (b) is the classical inequality between mean and median, in a normed space; (c)
identifies a fixed point, not its attraction; (d) uses $\polar(rx^\top)=rx^\top/(\norm r\norm x)$
\cite[Lemma~H.1]{li2026implicit}. For $K=1$ the limit in (d) is classical (sign-error algorithms minimize a mean
absolute error), for $d=1$ it is the geometric median of the $y_i/x_i$ \cite{cardot2013efficient}, and for
coordinatewise sign descent at batch size one non-convergence for every step-size sequence is
proved in \cite{karimireddy2019error} (Theorem~I) when all data points share a sign pattern; for $K,d\ge2$ (d) was not
found.

\begin{proposition}[size of the offset for minibatches; formal]\label{prop:offset}
Let $\xi_b$ be the mean of $b$ i.i.d.\ centred matrices with i.i.d.\ entries of variance $\sigma^2$ and third
cumulant $\kappa_3$, with a non-lattice law. A one-term Edgeworth expansion gives, for the minimizer $G_b^\dagger$ of
$\psi_{\xi_b}$,
\begin{equation}\label{eq:offset}
G_b^\dagger=\frac{\kappa_3}{6\sigma^2b}\cdot\frac{a_{K,d}}{c_{K,d}}\;\ones\ones^\top+o\Big(\frac1b\Big),\qquad
a_{K,d}=-\E\big[\polar(Z)_{11}(Z_{11}^3-3Z_{11})\big],\quad c_{K,d}=\frac{\E\nnuc Z}{Kd},
\end{equation}
with $Z$ a standard Gaussian $K\times d$ matrix: in every entry, the scalar offset $\kappa_3/(6\sigma^2b)$ of the
median of a mean (Cornish--Fisher \cite{cornish1938moments}) times $a_{K,d}/c_{K,d}$.
\end{proposition}
The expansion is formal (an Edgeworth expansion against a discontinuous map; the linear response $c_{K,d}$ is in
\cite{xie2026muon,chen2026when}), but its predictions are parameter-free and checked below. What matters is the
skewness of the \emph{gradient} noise: with symmetric inputs independent of the residuals the offset is zero, and an
intercept suffices to create one. The constant $a_{K,d}/c_{K,d}$, symmetric in $(K,d)$, is numerically of order $1/\max(K,d)$:
it equals $3/(d+2)$ for $K=1$, and $\max(K,d)\,a_{K,d}/c_{K,d}$ lies between $1.4$ and $2.8$ for ten sizes with
$K\le d$ (quadrature for $2\times3$, Monte Carlo up to $32\times32$ and $16\times64$, relative standard errors below
$3\%$): per entry, the offset of the polar step is smaller than that of sign descent by a factor of order
$\max(K,d)$. For Muon itself, in dimension one with the Newton--Schulz map and momentum, \cite{do2026muon} prove
non-convergence for bounded data with a nonzero third central moment, for almost all batch sizes (Theorem~1.3), and
show that every limit is a zero of the mean field (Proposition~6.2).

\paragraph{Numerical checks.} For $2\times3$ matrices with entries $\mathrm{Exp}(1)-1$,
\eqref{eq:offset} predicts $b\,\|G_b^\dagger\|_*\to0.464$; the zero of the mean field computed from $1.5\cdot10^6$
exactly sampled batch means gives $0.450$--$0.498$ for $b=2,\dots,256$, the extremes at $b=2$ and $b=256$, where
higher-order terms and the Monte-Carlo error of the root, respectively, are largest. On the finite sum of
Figure~\ref{fig:theory} ($\nF{W_{\mathrm{LAD}}-W_{\mathrm{LS}}}=0.170$), averaged batch-one polar steps with decreasing lengths
end $2.6\cdot10^{-3}$ from $W_{\mathrm{LAD}}$ and $0.169$ from $W_{\mathrm{LS}}$, the SAGA estimator with the same
lengths $2.1\cdot10^{-4}$ from $W_{\mathrm{LS}}$; for $b\le32$ the runs end within $5\%$ of the zero of the mean field
(exact for $b=1$, Monte Carlo otherwise, two estimates within $5\%$), and the SAGA estimator with the nuclear-scaled
length at rounding level (gradient norm below $2.4\cdot10^{-15}$).

\section{Discussion}\label{sec:discussion}

\subsection{From SAGA to Muon: what transferred, and why}\label{sec:link}

A Muon-type step is a change of geometry and tests both halves of Section~\ref{sec:classical}: on ridge
logistic and softmax models Theorem~\ref{thm:nfl} extends the worst-case half \finding{C5} to every norm, so
orthogonalization can only buy a property of trajectories, the other half, and Euclidean devices reproduce most of it
\finding{M3}. It is also nonlinear in the estimator, whereas SAGA rests on two
properties of its estimator: it is unbiased, which in a linear step makes the mean step the gradient step, and its
error vanishes at the minimizer, which removes the variance term. Through the polar map the first property is lost,
$\E\,\polar(G+\xi)\ne\polar(G)$: the mean step is a mirror subgradient step on a smoothed nuclear norm, minimized where
the gradient is the negative of a median of the noise, not of its mean (Theorem~\ref{thm:smd}). The second survives: at the minimizer,
with its memory there too, the estimator is exactly zero (Proposition~\ref{prop:comp}(c)). Variance reduction
therefore acts on the fixed point, which loses its offset, and not only on a variance term in a rate
(Figure~\ref{fig:theory}).
Table~\ref{tab:link} lists what was carried over. The link is one of tools, questions and this mechanism, not a
reduction: when the polar factor is recomputed at every step the proof of SAGA does not apply and we have no
convergence proof, and the results on noise without a finite sum (Theorem~\ref{thm:smd},
Proposition~\ref{prop:offset}) owe nothing to SAGA but the question.

\begin{table}[!htb]
\caption{From SAGA to Muon: what was carried over, and with which outcome.}\label{tab:link}
\centering\footnotesize
\renewcommand{\arraystretch}{1.1}
\begin{tabularx}{\textwidth}{@{}>{\raggedright\arraybackslash}p{0.33\textwidth}YY@{}}
\toprule
In the SAGA study & Counterpart for Muon-type steps & Outcome \\
\midrule
\emph{Tool.} One-step expectations computed exactly by enumerating the $n$ indices. &
The same enumeration, although the step is nonlinear in the estimator. &
Exact mean steps: the mechanism of \finding{M2} and the zeros of the mean field at batch one
(Figure~\ref{fig:theory}). \\
\emph{Mechanism.} An estimator that is unbiased and whose error vanishes at the minimizer. &
Unbiasedness is lost through the polar map; the vanishing error is kept. &
Variance reduction removes an offset (Theorem~\ref{thm:smd}, Proposition~\ref{prop:comp}); the step must vanish too
\finding{M1}. \\
\emph{Impossibility.} No fixed metric improves the logistic worst case \finding{C5}. &
No norm does, when the loss curvature vanishes (Theorem~\ref{thm:nfl}). &
On softmax, methods can only be compared on trajectories \finding{M3}. \\
\emph{Proof.} Proposition~5.9 in a fixed metric \finding{C5}. &
A polar metric computed once \finding{M4}. &
A certified rate; none when the polar factor is recomputed. \\
\emph{Controls.} SAGA in the covariance metric \finding{C5}; the memory ceiling \finding{C4}. &
Baseline and ceiling of the comparison. &
The covariance metric is the best method on least squares; no geometry helps where memory limits \finding{M3}. \\
\emph{Acceleration.} Negative momentum \finding{C6}. &
A spectral step inside loopless Katyusha. &
No gain over the Euclidean method. \\
\bottomrule
\end{tabularx}
\end{table}

\subsection{Relation to prior work}\label{sec:prior}

Beyond Table~\ref{tab:sota}: for the experiments, finite-sum
estimators in fixed-length polar steps \cite{li2026implicit}, in general norms \cite{allenzhu2018katyusha} and with a
preconditioner computed once \cite{yang2014data} are published; a SAGA table in the nuclear-scaled polar step at a
constant step size, the reading \finding{M2} and a frozen metric taken from the polar decomposition of one gradient
were not found by the check of 12 November 2026. Two topics of the theory note are left out because they are largely
published: the full-batch polar dynamics on least squares (exact line search with rate $1-1/\kop$ and a slow tail
\cite{shulgin2026quadratic}; fixed-length steps and the flow in the commuting case \cite{vasudeva2025muon}) and the
regime of low signal-to-noise ratio \cite{xie2026muon,chen2026when,paquette2026phases}. The exact facts that the note adds on the
former are not used here.

\subsection{Limitations}

Everything is for strongly convex objectives, on synthetic data of moderate size; the shares of \finding{C2} depend on
the order in which the inequalities are undone; \finding{M3} rests on one to three draws of the data per instance.
Theorem~\ref{thm:smd} assumes additive noise and no momentum (for minibatches of a finite sum the law of the noise
depends on the iterate; only batch one on least squares is treated); Proposition~\ref{prop:offset} is formal, and the
scaling of its constant in $1/\max(K,d)$ is numerical (exact for $\min(K,d)=1$).

\subsection{Conclusion}

The textbook proof of SAGA has the right worst-case order but, measured along trajectories, loses where it replaces
the curvature met along the run by global constants, and no fixed change of geometry, Euclidean or spectral, repairs
the worst case on ridge logistic and softmax models. For Muon-type steps, finite-sum variance reduction matters less
for speed, where Euclidean devices do as well, than for precision: with additive noise and decreasing lengths,
stochastic polar steps converge to a point where the gradient is minus a nuclear median of the noise, whereas a
finite-sum estimator makes the minimizer the only fixed point. Open questions are convergence proofs when the polar factor is recomputed at every step, momentum,
iterate-dependent minibatch noise, and real data.

\paragraph{Reproducibility and use of AI tools.} All numbers come from scripts in \nolinkurl{experiments/},
\nolinkurl{muon/experiments/} and \nolinkurl{muon/theory_checks/}. A large language model (Claude, Anthropic) was used
under the authors' supervision to derive and check the statements, write the scripts and run adversarial reviews of
the mathematics and the literature (theorem numbers of preprints are those read then).

{\footnotesize\setlength{\columnsep}{1.6em}
\begin{multicols}{2}[\section*{References}]
\sloppy
\bibliographystyle{abbrv}
\bibliography{report_spectral_muon}
\end{multicols}
}

\end{document}
